\chapter{Complexity — Machine}
%%% generated by the blueprint pipeline from TCSlib.Complexity.SpaceComplexity.ExpBFS.Machine
\section{Overview}

Declarations of the Lean module \texttt{TCSlib.Complexity.SpaceComplexity.ExpBFS.Machine}.

\section{Declarations}

\begin{definition}[First register index]
\lean{Complexity.LogProg.ExpBFS.rN}
\label{Complexity.LogProg.ExpBFS.rN}
\leanok
The first register index is the element $0$ of the 17-element index set
$\{0,1,\ldots,16\}$.
\end{definition}

\begin{definition}[Second register index]
\lean{Complexity.LogProg.ExpBFS.rLg}
\label{Complexity.LogProg.ExpBFS.rLg}
\leanok
The second register index is the element $1$ of the 17-element index set
$\{0,1,\ldots,16\}$.
\end{definition}

\begin{definition}[Third register index]
\lean{Complexity.LogProg.ExpBFS.rs}
\label{Complexity.LogProg.ExpBFS.rs}
\leanok
The third register index is the element $2$ of the 17-element index set
$\{0,1,\ldots,16\}$.
\end{definition}

\begin{definition}[Fourth register index]
\lean{Complexity.LogProg.ExpBFS.rs1}
\label{Complexity.LogProg.ExpBFS.rs1}
\leanok
The fourth register index is the element $3$ of the 17-element index set
$\{0,1,\ldots,16\}$.
\end{definition}

\begin{definition}[Fifth register index]
\lean{Complexity.LogProg.ExpBFS.rX}
\label{Complexity.LogProg.ExpBFS.rX}
\leanok
The fifth register index is the element $4$ of the 17-element index set
$\{0,1,\ldots,16\}$.
\end{definition}

\begin{definition}[Read-write index]
\lean{Complexity.LogProg.ExpBFS.rw}
\label{Complexity.LogProg.ExpBFS.rw}
\leanok
Let $[17]=\{0,1,\ldots,16\}$. The read-write index is the element $5$ of $[17]$.
\end{definition}

\begin{definition}[Key index]
\lean{Complexity.LogProg.ExpBFS.rK}
\label{Complexity.LogProg.ExpBFS.rK}
\leanok
Let $[17]=\{0,1,\ldots,16\}$. The key index is the element $6$ of $[17]$.
\end{definition}

\begin{definition}[Bit-string index]
\lean{Complexity.LogProg.ExpBFS.rBs}
\label{Complexity.LogProg.ExpBFS.rBs}
\leanok
Let $[17]=\{0,1,\ldots,16\}$. The bit-string index is the element $7$ of $[17]$.
\end{definition}

\begin{definition}[U index]
\lean{Complexity.LogProg.ExpBFS.rU}
\label{Complexity.LogProg.ExpBFS.rU}
\leanok
Let $[17]=\{0,1,\ldots,16\}$. The $U$ index is the element $8$ of $[17]$.
\end{definition}

\begin{definition}[V index]
\lean{Complexity.LogProg.ExpBFS.rV}
\label{Complexity.LogProg.ExpBFS.rV}
\leanok
Let $[17]=\{0,1,\ldots,16\}$. The $V$ index is the element $9$ of $[17]$.
\end{definition}

\begin{definition}[T index]
\lean{Complexity.LogProg.ExpBFS.rT}
\label{Complexity.LogProg.ExpBFS.rT}
\leanok
Let $[17]=\{0,1,\ldots,16\}$. The $T$ index is the element $10$ of $[17]$.
\end{definition}

\begin{definition}[Second T index]
\lean{Complexity.LogProg.ExpBFS.rT2}
\label{Complexity.LogProg.ExpBFS.rT2}
\leanok
Let $[17]=\{0,1,\ldots,16\}$. The second $T$ index is the element $11$ of $[17]$.
\end{definition}

\begin{definition}[Z index]
\lean{Complexity.LogProg.ExpBFS.rZ}
\label{Complexity.LogProg.ExpBFS.rZ}
\leanok
Let $[17]=\{0,1,\ldots,16\}$. The $Z$ index is the element $12$ of $[17]$.
\end{definition}

\begin{definition}[First distinguished index]
\lean{Complexity.LogProg.ExpBFS.rC}
\label{Complexity.LogProg.ExpBFS.rC}
\leanok
The distinguished index $r_C$ is the element $13$ of the set of indices $	ext{Fin }17$,
namely the integers from $0$ through $16$.
\end{definition}

\begin{definition}[Second distinguished index]
\lean{Complexity.LogProg.ExpBFS.rch}
\label{Complexity.LogProg.ExpBFS.rch}
\leanok
The distinguished index $r_{ch}$ is the element $14$ of the set of indices $	ext{Fin
}17$, namely the integers from $0$ through $16$.
\end{definition}

\begin{definition}[Third distinguished index]
\lean{Complexity.LogProg.ExpBFS.rc}
\label{Complexity.LogProg.ExpBFS.rc}
\leanok
The distinguished index $r_c$ is the element $15$ of the set of indices $	ext{Fin }17$,
namely the integers from $0$ through $16$.
\end{definition}

\begin{definition}[Fourth distinguished index]
\lean{Complexity.LogProg.ExpBFS.rc1}
\label{Complexity.LogProg.ExpBFS.rc1}
\leanok
The distinguished index $r_{c1}$ is the element $16$ of the set of indices $	ext{Fin
}17$, namely the integers from $0$ through $16$.
\end{definition}

\begin{definition}[Exponential BFS machine state]
\lean{Complexity.LogProg.ExpBFS.St}
\label{Complexity.LogProg.ExpBFS.St}
\leanok
A machine state is a tuple of 18 natural numbers
$(N,L_g,s,s_1,X,w,K,B_s,U,V,T,T_2,Z,C,ch,c,c_1)$, representing respectively the
input-length counter, logarithm bound, radius, next radius, width factor, code width,
code bound, memory-region base, source code, target code, outward-walk counter,
return-walk counter, copy-and-shift scratch register, shift counter, relaxation-round
change flag, marked-code count, and first-search count.
\end{definition}

\begin{definition}[State counter vector]
\lean{Complexity.LogProg.ExpBFS.St.toV}
\label{Complexity.LogProg.ExpBFS.St.toV}
\leanok
\uses{Complexity.LogProg.ExpBFS.St}
For a state $st$, its counter vector is the function from the 17 indices $0,\ldots,16$
to the natural numbers whose successive entries are
$N,Lg,s,s1,X,w,K,Bs,U,V,T,T2,Z,C,ch,c,c1$.
\end{definition}

\begin{definition}[Zero state]
\lean{Complexity.LogProg.ExpBFS.St.zero}
\label{Complexity.LogProg.ExpBFS.St.zero}
\leanok
\uses{Complexity.LogProg.ExpBFS.St}
The zero state is the state in which all 17 counters are zero.
\end{definition}

\begin{lemma}[Zero-state counter vector]
\lean{Complexity.LogProg.ExpBFS.St.toV_zero}
\label{Complexity.LogProg.ExpBFS.St.toV_zero}
\leanok
\uses{Complexity.LogProg.ExpBFS.St.toV, Complexity.LogProg.ExpBFS.St.zero}
\difficulty{4}
The counter vector of the exponential BFS state whose 17 counters are all zero is the
constant function on the 17 indices with value zero.
\end{lemma}

\begin{lemma}[Input-length coordinate]
\lean{Complexity.LogProg.ExpBFS.St.toV_N}
\label{Complexity.LogProg.ExpBFS.St.toV_N}
\leanok
\uses{Complexity.LogProg.ExpBFS.St, Complexity.LogProg.ExpBFS.St.toV, Complexity.LogProg.ExpBFS.rN}
\difficulty{1}
For an exponential breadth-first-search state $S$, the coordinate at index $0$ of its
17-entry counter vector is its input-length counter $N$; that is, this coordinate equals
$S.N$.
\end{lemma}

\begin{lemma}[Logarithm-bound vector coordinate]
\lean{Complexity.LogProg.ExpBFS.St.toV_Lg}
\label{Complexity.LogProg.ExpBFS.St.toV_Lg}
\leanok
\uses{Complexity.LogProg.ExpBFS.St, Complexity.LogProg.ExpBFS.St.toV, Complexity.LogProg.ExpBFS.rLg}
\difficulty{1}
For every exponential breadth-first search state $st$, the coordinate at index $1$ of
its counter vector is the logarithm-bound counter $Lg$.
\end{lemma}

\begin{lemma}[Radius coordinate]
\lean{Complexity.LogProg.ExpBFS.St.toV_s}
\label{Complexity.LogProg.ExpBFS.St.toV_s}
\leanok
\uses{Complexity.LogProg.ExpBFS.St, Complexity.LogProg.ExpBFS.St.toV, Complexity.LogProg.ExpBFS.rs}
\difficulty{1}
For any exponential BFS machine state $st$, the coordinate at index $2$ of its counter
vector is its radius $s$.
\end{lemma}

\begin{lemma}[Next-radius coordinate]
\lean{Complexity.LogProg.ExpBFS.St.toV_s1}
\label{Complexity.LogProg.ExpBFS.St.toV_s1}
\leanok
\uses{Complexity.LogProg.ExpBFS.St, Complexity.LogProg.ExpBFS.St.toV, Complexity.LogProg.ExpBFS.rs1}
\difficulty{1}
For any exponential BFS machine state $st$, the coordinate at index $3$ of its counter
vector is its next radius $s_1$.
\end{lemma}

\begin{lemma}[Width-factor coordinate]
\lean{Complexity.LogProg.ExpBFS.St.toV_X}
\label{Complexity.LogProg.ExpBFS.St.toV_X}
\leanok
\uses{Complexity.LogProg.ExpBFS.St, Complexity.LogProg.ExpBFS.St.toV, Complexity.LogProg.ExpBFS.rX}
\difficulty{1}
For an exponential breadth-first-search state $S$, the coordinate at index $4$ of its
17-entry counter vector is its width factor $X$; that is, this coordinate equals $S.X$.
\end{lemma}

\begin{lemma}[Code-width coordinate]
\lean{Complexity.LogProg.ExpBFS.St.toV_w}
\label{Complexity.LogProg.ExpBFS.St.toV_w}
\leanok
\uses{Complexity.LogProg.ExpBFS.St, Complexity.LogProg.ExpBFS.St.toV, Complexity.LogProg.ExpBFS.rw}
\difficulty{1}
For any exponential BFS machine state $st$, the coordinate at index $5$ of its counter
vector is its code width $w$.
\end{lemma}

\begin{lemma}[Code-bound vector coordinate]
\lean{Complexity.LogProg.ExpBFS.St.toV_K}
\label{Complexity.LogProg.ExpBFS.St.toV_K}
\leanok
\uses{Complexity.LogProg.ExpBFS.St, Complexity.LogProg.ExpBFS.St.toV, Complexity.LogProg.ExpBFS.rK}
\difficulty{1}
For every exponential breadth-first search state $st$, the coordinate at index $6$ of
its counter vector is the code-bound counter $K$.
\end{lemma}

\begin{lemma}[Memory-base vector coordinate]
\lean{Complexity.LogProg.ExpBFS.St.toV_Bs}
\label{Complexity.LogProg.ExpBFS.St.toV_Bs}
\leanok
\uses{Complexity.LogProg.ExpBFS.St, Complexity.LogProg.ExpBFS.St.toV, Complexity.LogProg.ExpBFS.rBs}
\difficulty{1}
For every exponential breadth-first search state $st$, the coordinate at index $7$ of
its counter vector is the memory-region base counter $B_s$.
\end{lemma}

\begin{lemma}[Source-code coordinate]
\lean{Complexity.LogProg.ExpBFS.St.toV_U}
\label{Complexity.LogProg.ExpBFS.St.toV_U}
\leanok
\uses{Complexity.LogProg.ExpBFS.St, Complexity.LogProg.ExpBFS.St.toV, Complexity.LogProg.ExpBFS.rU}
\difficulty{1}
For an exponential breadth-first-search state $S$, the coordinate at index $8$ of its
17-entry counter vector is its source code $U$; that is, this coordinate equals $S.U$.
\end{lemma}

\begin{lemma}[Target-code coordinate]
\lean{Complexity.LogProg.ExpBFS.St.toV_V}
\label{Complexity.LogProg.ExpBFS.St.toV_V}
\leanok
\uses{Complexity.LogProg.ExpBFS.St, Complexity.LogProg.ExpBFS.St.toV, Complexity.LogProg.ExpBFS.rV}
\difficulty{1}
For an exponential breadth-first-search state $S$, the coordinate at index $9$ of its
17-entry counter vector is its target code $V$; that is, this coordinate equals $S.V$.
\end{lemma}

\begin{lemma}[Outward-walk counter coordinate]
\lean{Complexity.LogProg.ExpBFS.St.toV_T}
\label{Complexity.LogProg.ExpBFS.St.toV_T}
\leanok
\uses{Complexity.LogProg.ExpBFS.St, Complexity.LogProg.ExpBFS.St.toV, Complexity.LogProg.ExpBFS.rT}
\difficulty{1}
For an exponential breadth-first-search state $S$, the coordinate at index $10$ of its
17-entry counter vector is its outward-walk counter $T$; that is, this coordinate equals
$S.T$.
\end{lemma}

\begin{lemma}[Return-walk counter coordinate]
\lean{Complexity.LogProg.ExpBFS.St.toV_T2}
\label{Complexity.LogProg.ExpBFS.St.toV_T2}
\leanok
\uses{Complexity.LogProg.ExpBFS.St, Complexity.LogProg.ExpBFS.St.toV, Complexity.LogProg.ExpBFS.rT2}
\difficulty{1}
For an exponential breadth-first-search state $S$, the coordinate at index $11$ of its
17-entry counter vector is its return-walk counter $T_2$; that is, this coordinate
equals $S.T_2$.
\end{lemma}

\begin{lemma}[Scratch-register coordinate]
\lean{Complexity.LogProg.ExpBFS.St.toV_Z}
\label{Complexity.LogProg.ExpBFS.St.toV_Z}
\leanok
\uses{Complexity.LogProg.ExpBFS.St, Complexity.LogProg.ExpBFS.St.toV, Complexity.LogProg.ExpBFS.rZ}
\difficulty{1}
For an exponential breadth-first-search state $S$, the coordinate at index $12$ of its
17-entry counter vector is its copy-and-shift scratch register $Z$; that is, this
coordinate equals $S.Z$.
\end{lemma}

\begin{lemma}[Shift-counter vector coordinate]
\lean{Complexity.LogProg.ExpBFS.St.toV_C}
\label{Complexity.LogProg.ExpBFS.St.toV_C}
\leanok
\uses{Complexity.LogProg.ExpBFS.St, Complexity.LogProg.ExpBFS.St.toV, Complexity.LogProg.ExpBFS.rC}
\difficulty{1}
For every exponential breadth-first search state $st$, the coordinate at index $13$ of
its counter vector is the shift counter $C$.
\end{lemma}

\begin{lemma}[Change-flag coordinate]
\lean{Complexity.LogProg.ExpBFS.St.toV_ch}
\label{Complexity.LogProg.ExpBFS.St.toV_ch}
\leanok
\uses{Complexity.LogProg.ExpBFS.St, Complexity.LogProg.ExpBFS.St.toV, Complexity.LogProg.ExpBFS.rch}
\difficulty{1}
For any exponential BFS machine state $st$, the coordinate at index $14$ of its counter
vector is its change flag $ch$.
\end{lemma}

\begin{lemma}[Marked-code count coordinate]
\lean{Complexity.LogProg.ExpBFS.St.toV_c}
\label{Complexity.LogProg.ExpBFS.St.toV_c}
\leanok
\uses{Complexity.LogProg.ExpBFS.St, Complexity.LogProg.ExpBFS.St.toV, Complexity.LogProg.ExpBFS.rc}
\difficulty{1}
For any exponential BFS machine state $st$, the coordinate at index $15$ of its counter
vector is its marked-code count $c$.
\end{lemma}

\begin{lemma}[First-search count coordinate]
\lean{Complexity.LogProg.ExpBFS.St.toV_c1}
\label{Complexity.LogProg.ExpBFS.St.toV_c1}
\leanok
\uses{Complexity.LogProg.ExpBFS.St, Complexity.LogProg.ExpBFS.St.toV, Complexity.LogProg.ExpBFS.rc1}
\difficulty{1}
For any exponential BFS machine state $st$, the coordinate at index $16$ of its counter
vector is its first-search count $c_1$.
\end{lemma}

\begin{lemma}[Updating the input length counter]
\lean{Complexity.LogProg.ExpBFS.St.upd_N}
\label{Complexity.LogProg.ExpBFS.St.upd_N}
\leanok
\uses{Complexity.LogProg.ExpBFS.St, Complexity.LogProg.ExpBFS.St.toV, Complexity.LogProg.ExpBFS.rN}
\difficulty{4}
Let a machine state have natural-number registers
$(N,L_g,s,s_1,X,w,K,B_s,U,V,T,T_2,Z,C,ch,c,c_1)$, and let its register vector list these
values in that order, indexed from $0$ to $16$. Replacing coordinate $0$ of the vector
by a natural number $k$ gives the register vector of the state obtained by setting $N=k$
and leaving the other registers unchanged.
\end{lemma}

\begin{lemma}[Updating the logarithm bound]
\lean{Complexity.LogProg.ExpBFS.St.upd_Lg}
\label{Complexity.LogProg.ExpBFS.St.upd_Lg}
\leanok
\uses{Complexity.LogProg.ExpBFS.St, Complexity.LogProg.ExpBFS.St.toV, Complexity.LogProg.ExpBFS.rLg}
\difficulty{4}
Let a machine state have natural-number registers
$(N,L_g,s,s_1,X,w,K,B_s,U,V,T,T_2,Z,C,ch,c,c_1)$, and let its register vector list these
values in that order, indexed from $0$ to $16$. Replacing coordinate $1$ of the vector
by a natural number $k$ gives the register vector of the state obtained by setting
$L_g=k$ and leaving the other registers unchanged.
\end{lemma}

\begin{lemma}[Updating the Radius Coordinate]
\lean{Complexity.LogProg.ExpBFS.St.upd_s}
\label{Complexity.LogProg.ExpBFS.St.upd_s}
\leanok
\uses{Complexity.LogProg.ExpBFS.St, Complexity.LogProg.ExpBFS.St.toV, Complexity.LogProg.ExpBFS.rs}
\difficulty{4}
For any machine state $st$ and natural number $k$, replacing coordinate $2$ of its
17-entry counter vector by $k$ gives the counter vector of the state obtained from $st$
by setting its radius $s$ to $k$.
\end{lemma}

\begin{lemma}[Updating the Next-Radius Coordinate]
\lean{Complexity.LogProg.ExpBFS.St.upd_s1}
\label{Complexity.LogProg.ExpBFS.St.upd_s1}
\leanok
\uses{Complexity.LogProg.ExpBFS.St, Complexity.LogProg.ExpBFS.St.toV, Complexity.LogProg.ExpBFS.rs1}
\difficulty{4}
For any machine state $st$ and natural number $k$, replacing coordinate $3$ of its
17-entry counter vector by $k$ gives the counter vector of the state obtained from $st$
by setting its next radius $s_1$ to $k$.
\end{lemma}

\begin{lemma}[Updating the width factor]
\lean{Complexity.LogProg.ExpBFS.St.upd_X}
\label{Complexity.LogProg.ExpBFS.St.upd_X}
\leanok
\uses{Complexity.LogProg.ExpBFS.St, Complexity.LogProg.ExpBFS.St.toV, Complexity.LogProg.ExpBFS.rX}
\difficulty{4}
For any machine state $st$ and natural number $k$, replacing coordinate $4$ of its
register vector by $k$ yields the register vector of the state obtained from $st$ by
setting its width factor $X$ to $k$.
\end{lemma}

\begin{lemma}[Updating the Code-Width Coordinate]
\lean{Complexity.LogProg.ExpBFS.St.upd_w}
\label{Complexity.LogProg.ExpBFS.St.upd_w}
\leanok
\uses{Complexity.LogProg.ExpBFS.St, Complexity.LogProg.ExpBFS.St.toV, Complexity.LogProg.ExpBFS.rw}
\difficulty{4}
For any machine state $st$ and natural number $k$, replacing coordinate $5$ of its
17-entry counter vector by $k$ gives the counter vector of the state obtained from $st$
by setting its code width $w$ to $k$.
\end{lemma}

\begin{lemma}[Updating the code bound]
\lean{Complexity.LogProg.ExpBFS.St.upd_K}
\label{Complexity.LogProg.ExpBFS.St.upd_K}
\leanok
\uses{Complexity.LogProg.ExpBFS.St, Complexity.LogProg.ExpBFS.St.toV, Complexity.LogProg.ExpBFS.rK}
\difficulty{4}
Let a machine state have natural-number registers
$(N,L_g,s,s_1,X,w,K,B_s,U,V,T,T_2,Z,C,ch,c,c_1)$, and let its register vector list these
values in that order, indexed from $0$ to $16$. Replacing coordinate $6$ of the vector
by a natural number $k$ gives the register vector of the state obtained by setting $K=k$
and leaving the other registers unchanged.
\end{lemma}

\begin{lemma}[Updating the memory base]
\lean{Complexity.LogProg.ExpBFS.St.upd_Bs}
\label{Complexity.LogProg.ExpBFS.St.upd_Bs}
\leanok
\uses{Complexity.LogProg.ExpBFS.St, Complexity.LogProg.ExpBFS.St.toV, Complexity.LogProg.ExpBFS.rBs}
\difficulty{4}
Let a machine state have natural-number registers
$(N,L_g,s,s_1,X,w,K,B_s,U,V,T,T_2,Z,C,ch,c,c_1)$, and let its register vector list these
values in that order, indexed from $0$ to $16$. Replacing coordinate $7$ of the vector
by a natural number $k$ gives the register vector of the state obtained by setting
$B_s=k$ and leaving the other registers unchanged.
\end{lemma}

\begin{lemma}[Updating the source code register]
\lean{Complexity.LogProg.ExpBFS.St.upd_U}
\label{Complexity.LogProg.ExpBFS.St.upd_U}
\leanok
\uses{Complexity.LogProg.ExpBFS.St, Complexity.LogProg.ExpBFS.St.toV, Complexity.LogProg.ExpBFS.rU}
\difficulty{4}
For any machine state $st$ and natural number $k$, replacing coordinate $8$ of its
register vector by $k$ yields the register vector of the state obtained from $st$ by
setting its source code register $U$ to $k$.
\end{lemma}

\begin{lemma}[Updating the target code register]
\lean{Complexity.LogProg.ExpBFS.St.upd_V}
\label{Complexity.LogProg.ExpBFS.St.upd_V}
\leanok
\uses{Complexity.LogProg.ExpBFS.St, Complexity.LogProg.ExpBFS.St.toV, Complexity.LogProg.ExpBFS.rV}
\difficulty{4}
For any machine state $st$ and natural number $k$, replacing coordinate $9$ of its
register vector by $k$ yields the register vector of the state obtained from $st$ by
setting its target code register $V$ to $k$.
\end{lemma}

\begin{lemma}[Updating the outward walk counter]
\lean{Complexity.LogProg.ExpBFS.St.upd_T}
\label{Complexity.LogProg.ExpBFS.St.upd_T}
\leanok
\uses{Complexity.LogProg.ExpBFS.St, Complexity.LogProg.ExpBFS.St.toV, Complexity.LogProg.ExpBFS.rT}
\difficulty{4}
Let a machine state have natural-number registers
$(N,L_g,s,s_1,X,w,K,B_s,U,V,T,T_2,Z,C,ch,c,c_1)$, and let its register vector list these
values in that order, indexed from $0$ to $16$. Replacing coordinate $10$ of the vector
by a natural number $k$ gives the register vector of the state obtained by setting $T=k$
and leaving the other registers unchanged.
\end{lemma}

\begin{lemma}[Updating the return walk counter]
\lean{Complexity.LogProg.ExpBFS.St.upd_T2}
\label{Complexity.LogProg.ExpBFS.St.upd_T2}
\leanok
\uses{Complexity.LogProg.ExpBFS.St, Complexity.LogProg.ExpBFS.St.toV, Complexity.LogProg.ExpBFS.rT2}
\difficulty{4}
Let a machine state have natural-number registers
$(N,L_g,s,s_1,X,w,K,B_s,U,V,T,T_2,Z,C,ch,c,c_1)$, and let its register vector list these
values in that order, indexed from $0$ to $16$. Replacing coordinate $11$ of the vector
by a natural number $k$ gives the register vector of the state obtained by setting
$T_2=k$ and leaving the other registers unchanged.
\end{lemma}

\begin{lemma}[Updating the scratch register]
\lean{Complexity.LogProg.ExpBFS.St.upd_Z}
\label{Complexity.LogProg.ExpBFS.St.upd_Z}
\leanok
\uses{Complexity.LogProg.ExpBFS.St, Complexity.LogProg.ExpBFS.St.toV, Complexity.LogProg.ExpBFS.rZ}
\difficulty{3}
For any machine state $st$ and natural number $k$, replacing coordinate $12$ of its
register vector by $k$ yields the register vector of the state obtained from $st$ by
setting its scratch register $Z$ to $k$.
\end{lemma}

\begin{lemma}[Updating the shift counter]
\lean{Complexity.LogProg.ExpBFS.St.upd_C}
\label{Complexity.LogProg.ExpBFS.St.upd_C}
\leanok
\uses{Complexity.LogProg.ExpBFS.St, Complexity.LogProg.ExpBFS.St.toV, Complexity.LogProg.ExpBFS.rC}
\difficulty{3}
Let a machine state have natural-number registers
$(N,L_g,s,s_1,X,w,K,B_s,U,V,T,T_2,Z,C,ch,c,c_1)$, and let its register vector list these
values in that order, indexed from $0$ to $16$. Replacing coordinate $13$ of the vector
by a natural number $k$ gives the register vector of the state obtained by setting $C=k$
and leaving the other registers unchanged.
\end{lemma}

\begin{lemma}[Updating the change flag]
\lean{Complexity.LogProg.ExpBFS.St.upd_ch}
\label{Complexity.LogProg.ExpBFS.St.upd_ch}
\leanok
\uses{Complexity.LogProg.ExpBFS.St, Complexity.LogProg.ExpBFS.St.toV, Complexity.LogProg.ExpBFS.rch}
\difficulty{3}
For any machine state $st$ and natural number $k$, replacing coordinate $14$ of its
register vector by $k$ yields the register vector of the state obtained from $st$ by
setting its change flag $ch$ to $k$.
\end{lemma}

\begin{lemma}[Updating the marked-code count]
\lean{Complexity.LogProg.ExpBFS.St.upd_c}
\label{Complexity.LogProg.ExpBFS.St.upd_c}
\leanok
\uses{Complexity.LogProg.ExpBFS.St, Complexity.LogProg.ExpBFS.St.toV, Complexity.LogProg.ExpBFS.rc}
\difficulty{3}
For any machine state $st$ and natural number $k$, replacing coordinate $15$ of its
register vector by $k$ yields the register vector of the state obtained from $st$ by
setting its marked-code count $c$ to $k$.
\end{lemma}

\begin{lemma}[Updating the first-search count]
\lean{Complexity.LogProg.ExpBFS.St.upd_c1}
\label{Complexity.LogProg.ExpBFS.St.upd_c1}
\leanok
\uses{Complexity.LogProg.ExpBFS.St, Complexity.LogProg.ExpBFS.St.toV, Complexity.LogProg.ExpBFS.rc1}
\difficulty{3}
For any machine state $st$ and natural number $k$, replacing coordinate $16$ of its
register vector by $k$ yields the register vector of the state obtained from $st$ by
setting its first-search count $c_1$ to $k$.
\end{lemma}

\begin{definition}[Uniform state bound]
\lean{Complexity.LogProg.ExpBFS.St.Bd}
\label{Complexity.LogProg.ExpBFS.St.Bd}
\leanok
\uses{Complexity.LogProg.ExpBFS.St}
For a natural number $B$ and a state $st$, the state is bounded by $B$ when each of its
17 counters is at most $B$.
\end{definition}

\begin{lemma}[Bound on state-vector coordinates]
\lean{Complexity.LogProg.ExpBFS.St.Bd.toV}
\label{Complexity.LogProg.ExpBFS.St.Bd.toV}
\leanok
\uses{Complexity.LogProg.ExpBFS.St, Complexity.LogProg.ExpBFS.St.Bd, Complexity.LogProg.ExpBFS.St.toV}
\difficulty{4}
Let $B$ be a natural number and let $st$ be an exponential breadth-first search state
whose 17 counters are each at most $B$. Then every coordinate of the state’s 17-entry
counter vector is at most $B$.
\end{lemma}

\begin{definition}[Three search phases]
\lean{Complexity.LogProg.ExpBFS.Ph3}
\label{Complexity.LogProg.ExpBFS.Ph3}
\leanok
There are three phases: the first search, the second search, and the final scan.
\end{definition}

\begin{definition}[Phase radius]
\lean{Complexity.LogProg.ExpBFS.rad}
\label{Complexity.LogProg.ExpBFS.rad}
\leanok
\uses{Complexity.LogProg.ExpBFS.Ph3, Complexity.LogProg.ExpBFS.rs, Complexity.LogProg.ExpBFS.rs1}
Map the second search phase to radius $3$ and the first search and final scan phases to
radius $2$, with these radii in $\mathrm{Fin}(17)$.
\end{definition}

\begin{definition}[Boolean phase assignment]
\lean{Complexity.LogProg.ExpBFS.toPh}
\label{Complexity.LogProg.ExpBFS.toPh}
\leanok
\uses{Complexity.LogProg.ExpBFS.Ph3}
Assign the second search to the Boolean value true and the first search to false.
\end{definition}

\begin{definition}[Search sites]
\lean{Complexity.LogProg.ExpBFS.Site}
\label{Complexity.LogProg.ExpBFS.Site}
\leanok
A site is one of the following: marking an initial code, reading the mark of $U$ in a
relaxation round, reading the mark of $V$ in a relaxation round, marking $V$ in a
relaxation round, reading the mark of $U$ in the count, or reading the mark of $U$ in
the final scan. Each of the first five kinds of site carries a Boolean parameter; the
final-scan site carries none.
\end{definition}

\begin{definition}[Site index]
\lean{Complexity.LogProg.ExpBFS.Site.idx}
\label{Complexity.LogProg.ExpBFS.Site.idx}
\leanok
\uses{Complexity.LogProg.ExpBFS.Site, Complexity.LogProg.ExpBFS.rU, Complexity.LogProg.ExpBFS.rV}
Assign index $9$ to sites that read the mark of $V$ or mark $V$ in a relaxation round,
and index $8$ to every other site, with indices in $\mathrm{Fin}(17)$.
\end{definition}

\begin{definition}[Marked sites]
\lean{Complexity.LogProg.ExpBFS.Site.set}
\label{Complexity.LogProg.ExpBFS.Site.set}
\leanok
\uses{Complexity.LogProg.ExpBFS.Site}
A site is marked exactly when it marks an initial code or marks $V$ in a relaxation
round; all other sites are unmarked.
\end{definition}

\begin{definition}[Access phases]
\lean{Complexity.LogProg.ExpBFS.AcPh}
\label{Complexity.LogProg.ExpBFS.AcPh}
\leanok
An access phase is either an outward walk with a phase $j\in\{0,1,2,3\}$, a memory
operation, or a return walk carrying a bit $b\in\{\mathrm{false},\mathrm{true}\}$ and a
phase $j\in\{0,1,2\}$.
\end{definition}

\begin{definition}[Circuit space components]
\lean{Complexity.LogProg.ExpBFS.CpS}
\label{Complexity.LogProg.ExpBFS.CpS}
\leanok
\uses{Complexity.LogProg.ExpBFS.Ph3, Complexity.LogProg.ExpBFS.Site}
For each natural number $K_0$, the circuit-space components are: a radius into $X$ for
each search phase; a length parameter $Lg$ into $X$ for each phase; each of the $K_0$
additions of $X$ into $w$ for each phase; $w$ into the shift counter $C$ for each phase;
two components for placing $K$ into the base $Bs$, each indexed by a phase; access
components for placing the base $Bs$ and an index into $T$, each indexed by a site; an
iteration component placing $s$ into $s1$; and the component types described above.
\end{definition}

\begin{definition}[Program labels for the search procedure]
\lean{Complexity.LogProg.ExpBFS.Lb}
\label{Complexity.LogProg.ExpBFS.Lb}
\leanok
\uses{Complexity.LogProg.ExpBFS.AcPh, Complexity.LogProg.ExpBFS.CpS, Complexity.LogProg.ExpBFS.Ph3, Complexity.LogProg.ExpBFS.Site}
For a natural number $K_0$, let $\mathrm{Lb}(K_0)$ be the inductive type of labels for
the search procedure. Its labels represent setup and width operations, the first and
second searches and their relaxation rounds, counting and iteration steps, the final
scan and its accept or reject outcomes, copy-loop phases indexed by a copy site and a
phase in $\{0,\ldots,6\}$, move phases indexed by $\{0,1,2\}$, shift phases indexed by
one of three search stages and a step in $\{0,\ldots,8\}$, and memory-access phases
indexed by a site and an access phase; the copy sites include an addition index in
$\{0,\ldots,K_0-1\}$.
\end{definition}

\begin{definition}[Width-work exit labels]
\lean{Complexity.LogProg.ExpBFS.wkExit}
\label{Complexity.LogProg.ExpBFS.wkExit}
\leanok
\uses{Complexity.LogProg.ExpBFS.Lb, Complexity.LogProg.ExpBFS.Ph3}
For a natural number $K_0$, assign to each of the three stages—the first search, the
second search, and the final scan—its width-work exit label: respectively, the
initial-marking test with Boolean parameter false, the initial-marking test with Boolean
parameter true, and the final-scan test.
\end{definition}

\begin{definition}[Copy source register]
\lean{Complexity.LogProg.ExpBFS.cpSrc}
\label{Complexity.LogProg.ExpBFS.cpSrc}
\leanok
\uses{Complexity.LogProg.ExpBFS.CpS, Complexity.LogProg.ExpBFS.Site.idx, Complexity.LogProg.ExpBFS.rBs, Complexity.LogProg.ExpBFS.rK, Complexity.LogProg.ExpBFS.rLg, Complexity.LogProg.ExpBFS.rX, Complexity.LogProg.ExpBFS.rad, Complexity.LogProg.ExpBFS.rs, Complexity.LogProg.ExpBFS.rw}
For a natural number $K_0$, the source-register map assigns a register among $0$ through
$16$ to each copy site. It sends the radius copy to the radius register for its stage,
the logarithm copy to register $1$, each indexed addition copy to register $4$, the
width-counter copy to register $5$, both base-copy sites to register $6$, the
base-access copy to register $7$, the index-access copy to the register assigned to its
site (register $9$ for reading or marking $V$, and register $8$ otherwise), and the
iteration copy to register $2$.
\end{definition}

\begin{definition}[Copy destination register]
\lean{Complexity.LogProg.ExpBFS.cpDst}
\label{Complexity.LogProg.ExpBFS.cpDst}
\leanok
\uses{Complexity.LogProg.ExpBFS.CpS, Complexity.LogProg.ExpBFS.rBs, Complexity.LogProg.ExpBFS.rC, Complexity.LogProg.ExpBFS.rT, Complexity.LogProg.ExpBFS.rX, Complexity.LogProg.ExpBFS.rs1, Complexity.LogProg.ExpBFS.rw}
For a natural number $K_0$, the destination-register map assigns a register among $0$
through $16$ to each copy site. It sends the radius and logarithm copies to register
$4$, each indexed addition copy to register $5$, the width-counter copy to register
$13$, both base-copy sites to register $7$, both access copies to register $10$, and the
iteration copy to register $3$.
\end{definition}

\begin{definition}[Copy-loop exit labels]
\lean{Complexity.LogProg.ExpBFS.cpExit}
\label{Complexity.LogProg.ExpBFS.cpExit}
\leanok
\uses{Complexity.LogProg.ExpBFS.CpS, Complexity.LogProg.ExpBFS.Lb, Complexity.LogProg.ExpBFS.wkExit}
For a natural number $K_0$, assign an exit label to each copy site. A radius copy
proceeds to the logarithm copy at the same stage, and a logarithm copy exits to the
width increment at that stage. An indexed addition copy proceeds to the next addition
when its index plus one is less than $K_0$, and otherwise exits to the width clear at
that stage. A width-counter copy proceeds to the shift phase at step zero; a first base
copy exits to the width-work exit at the second-search stage and otherwise proceeds to
the second base copy; a second base copy exits to the width-work exit at its stage. A
base-access copy proceeds to the index-access copy at the same site, an index-access
copy proceeds to the outward memory-access phase at that site, and the iteration copy
exits to the increment of $s_1$.
\end{definition}

\begin{definition}[Memory-access exit labels]
\lean{Complexity.LogProg.ExpBFS.acExit}
\label{Complexity.LogProg.ExpBFS.acExit}
\leanok
\uses{Complexity.LogProg.ExpBFS.Lb, Complexity.LogProg.ExpBFS.Site}
For a natural number $K_0$, assign an exit label to each memory-access site and Boolean
outcome. At an initial-code marking site, either outcome proceeds to the next
initial-marking index with the same Boolean parameter. At a relaxation-round read of
$U$, true proceeds to the read of $V$ and false to the next $U$; at a read of $V$, true
proceeds to the next $V$ and false to the step-decider query. At a marking site for $V$,
either outcome proceeds to clear the change flag. At a count read of $U$, true proceeds
to increment the count and false to the next $U$. At the final-scan read of $U$, true
proceeds to the accepting-code query and false to the next $U$.
\end{definition}

\begin{definition}[Copy-loop instruction by phase]
\lean{Complexity.LogProg.ExpBFS.cpI}
\label{Complexity.LogProg.ExpBFS.cpI}
\leanok
\uses{Complexity.LogProg.ExpBFS.Lb, Complexity.LogProg.ExpBFS.rZ, Complexity.LogProg.Ins}
For a natural number $K_0$, let the source and destination be registers among $17$
registers, let $e$ be an exit label, let $m$ assign a label to each of seven phases, and
let $j$ be one of those phases. Define the instruction at phase $j$ as follows: at phase
$0$, test whether the source register is zero, branching to phase $4$ if it is and phase
$1$ otherwise; at phase $1$, decrement the source and proceed to phase $2$; at phase
$2$, increment the destination and proceed to phase $3$; at phase $3$, increment
register $0$ and proceed to phase $0$; at phase $4$, test whether register $0$ is zero,
branching to $e$ if it is and phase $5$ otherwise; at phase $5$, decrement register $0$
and proceed to phase $6$; and at phase $6$, increment the source and proceed to phase
$4$.
\end{definition}

\begin{definition}[Shift-loop instruction by phase]
\lean{Complexity.LogProg.ExpBFS.shI}
\label{Complexity.LogProg.ExpBFS.shI}
\leanok
\uses{Complexity.LogProg.ExpBFS.Lb, Complexity.LogProg.ExpBFS.rC, Complexity.LogProg.ExpBFS.rZ, Complexity.LogProg.Ins}
For a natural number $K_0$, given a register $r$ among $0$ through $16$, an exit label,
a function assigning a label to each of the nine shift-loop phases, and a phase
$j\in\{0,\ldots,8\}$, define the instruction as follows: phase $0$ tests whether
register $13$ is zero, branching to the exit if so and phase $1$ otherwise; phase $1$
decrements register $13$ and proceeds to phase $2; phase $2$ tests whether $r$ is zero,
branching to phase $6$ if so and phase $3$ otherwise; phase $3$ decrements $r$ and
proceeds to phase $4$; phases $4$ and $5$ each increment register $0$, proceeding
respectively to phase $5$ and back to phase $2$; phase $6$ tests whether register $0$ is
zero, branching to phase $0$ if so and phase $7$ otherwise; phase $7$ decrements
register $0$ and proceeds to phase $8; and phase $8$ increments $r$ and returns to phase
$6$.
\end{definition}

\begin{definition}[Counter increment loop]
\lean{Complexity.LogProg.ExpBFS.mvI}
\label{Complexity.LogProg.ExpBFS.mvI}
\leanok
\uses{Complexity.LogProg.ExpBFS.Lb, Complexity.LogProg.ExpBFS.rc, Complexity.LogProg.ExpBFS.rc1, Complexity.LogProg.Ins}
For a label $e$ and a phase $j\in\{0,1,2\}$, define a register instruction as follows:
at phase $0$, branch on whether register $r_c$ is zero, going to $e$ if it is and to
phase $1$ otherwise; at phase $1$, decrement $r_c$ and go to phase $2$; at phase $2$,
increment $r_{c_1}$ and return to phase $0$.
\end{definition}

\begin{definition}[Memory access routine]
\lean{Complexity.LogProg.ExpBFS.acI}
\label{Complexity.LogProg.ExpBFS.acI}
\leanok
\uses{Complexity.LogProg.ExpBFS.AcPh, Complexity.LogProg.ExpBFS.Lb, Complexity.LogProg.ExpBFS.Site, Complexity.LogProg.ExpBFS.Site.set, Complexity.LogProg.ExpBFS.acExit, Complexity.LogProg.ExpBFS.rT, Complexity.LogProg.ExpBFS.rT2, Complexity.LogProg.MIns}
For each memory site, define a memory instruction at each phase of an access routine.
The outward walk repeatedly tests register $r_T$: if it is zero, proceed to the memory
operation; otherwise decrement $r_T$, increment $r_{T2}$, and move the memory head right
before repeating. At the operation, set the current memory bit to $1$ at marking sites,
and otherwise branch on that bit; then return by repeatedly testing and decrementing
$r_{T2}$, moving left each time, and upon reaching zero continue to the site’s
prescribed exit label for the resulting bit.
\end{definition}

\begin{definition}[Register selected for clearing]
\lean{Complexity.LogProg.ExpBFS.clrReg}
\label{Complexity.LogProg.ExpBFS.clrReg}
\leanok
\uses{Complexity.LogProg.ExpBFS.rBs, Complexity.LogProg.ExpBFS.rK, Complexity.LogProg.ExpBFS.rU, Complexity.LogProg.ExpBFS.rV, Complexity.LogProg.ExpBFS.rc, Complexity.LogProg.ExpBFS.rc1, Complexity.LogProg.ExpBFS.rs1, Complexity.LogProg.ExpBFS.rw}
For each natural number $i$, assign a register among the eight registers used in the
cleanup loop: the assignments for $i=0,1,2,3,4,5,6$ are respectively the registers
$r_U,r_V,r_w,r_K,r_{Bs},r_c,r_{c_1}$, and every $i\ge 7$ is assigned $r_{s_1}$.
\end{definition}

\begin{definition}[Breadth-first search program]
\lean{Complexity.LogProg.ExpBFS.A}
\label{Complexity.LogProg.ExpBFS.A}
\leanok
\uses{Complexity.LogProg.ExpBFS.Lb, Complexity.LogProg.ExpBFS.acI, Complexity.LogProg.ExpBFS.clrReg, Complexity.LogProg.ExpBFS.cpDst, Complexity.LogProg.ExpBFS.cpExit, Complexity.LogProg.ExpBFS.cpI, Complexity.LogProg.ExpBFS.cpSrc, Complexity.LogProg.ExpBFS.mvI, Complexity.LogProg.ExpBFS.rK, Complexity.LogProg.ExpBFS.rLg, Complexity.LogProg.ExpBFS.rN, Complexity.LogProg.ExpBFS.rU, Complexity.LogProg.ExpBFS.rV, Complexity.LogProg.ExpBFS.rX, Complexity.LogProg.ExpBFS.rad, Complexity.LogProg.ExpBFS.rc, Complexity.LogProg.ExpBFS.rc1, Complexity.LogProg.ExpBFS.rch, Complexity.LogProg.ExpBFS.rs, Complexity.LogProg.ExpBFS.rs1, Complexity.LogProg.ExpBFS.shI, Complexity.LogProg.ExpBFS.toPh, Complexity.LogProg.MARM}
For each natural number $K_0$, define a program for a nondeterministic finite Turing
machine on binary inputs. Its labels implement the following stages: query the
input-length decider while counting the input length and computing its logarithmic
scale; initialize the search width; perform two searches using the initial-code and step
deciders, relaxing marks and counting marked codes; compare the counts; and then scan
for an accepting code, halting with acceptance or rejection. The program uses copy
routines, a counter-transfer routine, shifts of the search bound, and memory-access
routines to carry out these stages.
\end{definition}
